\chapter{Questionnaire items and details}
\label{AppendixB}

These are the items of the questionnaire required to access MediaBots. With the exception of the first item, which is in the form of an open-ended question, all other items are in the form of a 6-point Likert-scale question with options:

\begin{enumerate}
\item Strongly Disagree
\item Disagree
\item Slightly Disagree
\item Slightly Agree
\item Agree
\item Strongly Agree
\end{enumerate}

The first 5 items are aimed at measuring the participant's knowledge and intuition about how this technology works, as well as their expectations about image \emph{alignment} with the text prompts.  

\begin{enumerate}
\item Based on what you currently know, how will text-to-image technology affect your creative process? [Open question]
\item I understand how this technology interprets keywords to generate images.
\item I can imagine how this technology interprets keywords to generate images.
\item I expect this technology to interpret the keywords in a similar way as humans do
\item I expect generated pictures to be an accurate visual representation of the keywords
\end{enumerate}

The next 10 questions are attempting to measure the psychological construct of tolerance for ambiguity, adapted from \cite{Herman2010}:

\begin{enumerate}[resume]
\item I like to surround myself with things that are familiar to me.
\item The sooner we all share similar ideals the better.
\item A desirable job is one where what is to be done is always clear.
\item I prefer a life with few surprises.
\item What we are used to is preferable to what is unfamiliar.
\item I prefer settings where people share my values.
\item I like parties where the attendees are strangers.
\item I can enjoy being with people whose values are different from mine.
\item I would like to live in a foreign country for a while.
\item I can be comfortable with all kinds of people.
\end{enumerate}

